\documentclass{bmstu-sob}

\begin{document}

\tableofcontents

\bmstuIntroduction

Это проверочный текст класса \texttt{bmstu-sob}. Основной шрифт — Times New Roman, размер 14 пт, межстрочный интервал полуторный. Абзацный отступ составляет 1,25 см.

\section{Тестовая секция}
Текст основной секции. По требованиям отчёта новый раздел начинается с новой страницы.

\subsection{Тестовый подраздел}
Текст подраздела.

\subsubsection{Тестовый пункт}
Текст пункта.

\paragraph{Пункт с заголовком}
Этот уровень является run-in заголовком и не создаёт большой вертикальный пробел.

\begin{itemize}
  \item первый элемент;
  \item второй элемент.
\end{itemize}

\begin{enumerate}
  \item первый элемент;
  \item второй элемент.
\end{enumerate}

\begin{equation}
  A = \frac{B}{C}
  \label{eq:test}
\end{equation}

\begin{bmstuwherelist}
  \bmstuwhereline{$B$ — числитель;}
  \bmstuwhereline{$C$ — знаменатель.}
\end{bmstuwherelist}

Ссылка на формулу: \eqrefn{eq:test}.

\begin{figure}[htbp]
  \centering
  \fbox{\rule{0pt}{35mm}\rule{75mm}{0pt}}
  \caption{Проверочный рисунок}
  \label{fig:test}
\end{figure}

\begin{table}[htbp]
  \caption{Проверочная таблица}
  \label{tab:test}
  \centering
  \begin{tabular}{ll}
    Параметр & Значение \\
    A & 10 \\
    B & 20 \\
  \end{tabular}
  \bmstuTableNote{пример примечания к таблице.}
\end{table}

\bmstuConclusion

Проверочный текст заключения.

\bmstuReferencesTitle

\begin{bmstubibliography}
  \bmstubibitem{1}{Иванов И. И. Пример библиографического описания. Москва : Пример, 2024.}
\end{bmstubibliography}

\appendix
\bmstuAppendix{Проверочное приложение}

Текст приложения.

\section{Раздел приложения}
Текст раздела приложения.

\end{document}
